\documentclass[11pt,a4paper]{article}
\usepackage[utf8]{inputenc}
\usepackage[T1]{fontenc}
\usepackage{times}
\usepackage{amsmath,amssymb}
\usepackage{booktabs}
\usepackage{graphicx}
\usepackage{multirow}
\usepackage{array}
\usepackage{tabularx}
\usepackage[margin=1in]{geometry}
\usepackage{natbib}
\usepackage{hyperref}
\usepackage{xcolor}
\usepackage{caption}

\newcommand{\lzstar}{$l_{z}^{*}$}
\newcommand{\lzm}{$l_{zm}$}

\title{IRT-FIT: Detecting Benchmark Contamination in Large Language Models\\via Item Response Theory Person-Fit Statistics}

\author{Anonymous}
\date{}

\begin{document}
\maketitle

%% ============================================================
%% ABSTRACT
%% ============================================================
\begin{abstract}
Benchmark contamination---the inclusion of evaluation items in training data---undermines the reliability of Large Language Model (LLM) evaluation.
Existing detection methods require privileged access to training corpora, token-level probabilities, or model weights, limiting their applicability to retrospective audits.
We introduce \textsc{IRT-FIT}, a contamination detection framework operating at the \emph{minimal-information} level: only binary (correct/incorrect) response patterns are required.
Drawing on Item Response Theory (IRT) person-fit statistics developed for exam cheating detection, we characterize the operating envelope in which detection is effective.
Simulation and theoretical analysis establish that detection power exceeds 90\% on benchmarks with $\geq$10K items at 15\% contamination, but follows an inverted-U curve---a detection trough emerges at 40--70\% contamination where partial ability inflation masks aberrant patterns.
Controlled contamination experiments on Qwen3.5-9B and Gemma-4-31B confirm a measurable $\Delta l_{z}^{*}$ shift (Cohen's $d$=1.2 on MMLU at 15\%), with the inverted-U observed in real SFT data.
Comparison with ConStat and Min-K\%++ quantifies the cost of minimal information: a 13-point AUROC gap that constitutes a principled power--information tradeoff.
A leaderboard audit of 310 models from the Open LLM Leaderboard identifies 18 models flagged across $\geq$3 benchmarks, with contamination prevalence increasing from 12\% (2023~H1) to 25\% (2024~H1).
The quantified operating envelope---where IRT-FIT works, where it does not, and at what cost---is itself a contribution to the contamination detection literature.
\end{abstract}

%% ============================================================
%% 1. INTRODUCTION
%% ============================================================
\section{Introduction}

As LLM evaluation increasingly relies on standardized benchmarks such as MMLU, HellaSwag, and GSM8K, the integrity of these benchmarks becomes a critical concern.
Benchmark contamination---whether through deliberate inclusion of evaluation items in training data or incidental overlap in web-crawled corpora---inflates reported performance and distorts model comparisons \citep{zhou2023contamination,yang2023rethinking}.

All existing contamination detection methods require access beyond what a binary score matrix provides.
Membership Inference Attacks (MIA) such as Min-K\%++ \citep{shi2024detecting} require token-level log-probabilities.
ConStat \citep{dekoninck2024constat} requires live model access for rephrased retesting.
Time Travel probing \citep{oren2024proving} needs crafted prompts.
CoDeC \citep{shi2026codec} relies on in-context learning dynamics.
Critically, \citet{wang2025postrl} demonstrate that all MIA-based methods fail after reinforcement learning fine-tuning, as RL redistributes token-level probability distributions.

We observe that benchmark contamination is functionally identical to a well-studied problem in psychometrics: detecting examinees with item preknowledge.
Item Response Theory (IRT) person-fit statistics---specifically, the standardized log-likelihood statistic $l_{z}^{*}$ \citep{snijders2001}---have been used for decades to identify aberrant response patterns in educational testing.
These statistics require only binary responses, making them uniquely suited to the minimal-information operating point.

We introduce \textsc{IRT-FIT}, a principled framework for contamination detection that establishes and characterizes this operating point.
Rather than claiming superiority over existing methods, we precisely quantify what can be detected from binary responses alone, and at what cost.

Our contributions are:
\begin{enumerate}
    \item \textbf{Theoretical characterization.} We derive the detection power as a function of benchmark size ($n_\text{items}$) and contamination rate, revealing an inverted-U power curve with a detection trough at 40--70\% contamination. Theory and simulation agree with $R^2=0.85$.
    \item \textbf{Controlled calibration.} Controlled contamination experiments on two model families (Qwen3.5-9B, Gemma-4-31B) across five benchmarks establish ground-truth detection sensitivity, revealing an operating envelope of $\geq$15\% contamination on benchmarks with $\geq$5K items.
    \item \textbf{Power--information tradeoff.} Head-to-head comparison with ConStat and Min-K\%++ quantifies a 13-point AUROC gap as the cost of minimal information, establishing a detection frontier from binary responses to token probabilities.
    \item \textbf{Leaderboard audit.} Application to 310 Open LLM Leaderboard models identifies 18 multi-benchmark flagged models and a temporal contamination trend, demonstrating IRT-FIT's unique capability for retrospective audit of historical submissions.
\end{enumerate}

%% ============================================================
%% 2. BACKGROUND AND RELATED WORK
%% ============================================================
\section{Background}

\subsection{IRT and the 4PL Model}

Item Response Theory models the probability that examinee $j$ correctly answers item $i$ as a function of examinee ability $\theta_j$ and item parameters.
We adopt the four-parameter logistic (4PL) model:
\begin{equation}
    P(X_{ij}=1 \mid \theta_j) = c_i + (d_i - c_i) \cdot \frac{1}{1 + \exp(-a_i(\theta_j - b_i))}
\end{equation}
where $a_i$ is discrimination, $b_i$ is difficulty, $c_i$ is the guessing parameter, and $d_i$ is the feasibility (upper asymptote) parameter.
We use item parameters from PSN-IRT \citep{psnirt2025}, a large-scale IRT calibration of LLM benchmark items.

\subsection{Person-Fit Statistics}

The log-likelihood person-fit statistic for examinee $j$ is:
\begin{equation}
    l_0(\theta_j) = \sum_{i=1}^{n} \left[ x_{ij} \log P_{ij} + (1 - x_{ij}) \log(1 - P_{ij}) \right]
\end{equation}
The standardized variant $l_{z}^{*}$ \citep{snijders2001} accounts for the estimation error in $\hat{\theta}$:
\begin{equation}
    l_{z}^{*} = \frac{l_0 - E[l_0]}{\sqrt{\text{Var}(l_0)}}
\end{equation}
Under the null hypothesis (no aberrance), $l_{z}^{*} \sim \mathcal{N}(0,1)$.
Negative values indicate aberrant patterns (unexpected correct answers on hard items); positive values indicate over-consistent patterns (too predictable).

\subsection{Contamination Detection Methods}

We position \textsc{IRT-FIT} within a detection frontier organized by information access:
\begin{itemize}
    \item \textbf{Level 1} (binary responses): IRT-FIT (this work)
    \item \textbf{Level 2} (binary + retest): ConStat \citep{dekoninck2024constat}
    \item \textbf{Level 3} (token probabilities): Min-K\%++ \citep{shi2024detecting}
    \item \textbf{Level 4+} (weights/training data): n-gram overlap, internal probing
\end{itemize}

%% ============================================================
%% 3. METHOD
%% ============================================================
\section{The IRT-FIT Framework}

\subsection{Detection Logic}

When a model is contaminated with benchmark items through SFT, it acquires ``preknowledge'' of specific items.
This creates an aberrant response pattern: the model answers items correctly that its estimated ability $\hat{\theta}$ would not predict.
At low-to-moderate contamination, this produces $l_{z}^{*} \ll 0$ (unexpected correct answers).
At near-total contamination, the model appears uniformly high-performing, producing $l_{z}^{*} \gg 0$ (over-consistency).

\subsection{The Inverted-U Hypothesis}

We hypothesize---and subsequently verify---that detection power follows an inverted-U as a function of contamination rate:
\begin{itemize}
    \item \textbf{Low contamination} ($<$5\%): Signal drowned in sampling noise.
    \item \textbf{Moderate contamination} (10--25\%): Peak detection via negative $l_{z}^{*}$ (aberrance).
    \item \textbf{Transition zone} (40--70\%): Detection trough. Ability estimate $\hat{\theta}$ inflated by contamination, reducing apparent aberrance, while positive signal not yet established.
    \item \textbf{Near-total contamination} ($>$85\%): Detection via positive $l_{z}^{*}$ (over-consistency).
\end{itemize}

\subsection{Calibration via Controlled Contamination}

Real LLM responses exhibit universal IRT model misfit---all models produce $l_{z}^{*}$ values far from $\mathcal{N}(0,1)$, due to violations of local independence and unidimensionality.
This renders absolute $l_{z}^{*}$ thresholds unreliable.
We address this via controlled contamination: training clean and contaminated variants of the same model on identical data except for benchmark item inclusion, then measuring $\Delta l_{z}^{*}$ as the contamination-induced shift above natural misfit.

\subsection{Multiscale Aggregation ($l_{zm}$)}

Following \citet{hong2021multiscale}, we aggregate per-benchmark $l_{z}^{*}$ into a model-level score:
\begin{equation}
    l_{zm} = \frac{\sum_b w_b \cdot l_{z,b}^{*}}{\sqrt{\sum_b w_b^2}}, \quad w_b = \frac{1}{\text{SE}(l_{z,b}^{*})}
\end{equation}
This inverse-variance weighting naturally up-weights large benchmarks (low SE) and down-weights small ones.

%% ============================================================
%% 4. SIMULATION STUDIES
%% ============================================================
\section{Simulation Studies}

\subsection{Detection Power Simulation}
\label{sec:sim-power}

We simulate contamination by modifying IRT response vectors: for a given contamination rate $c$, a randomly selected fraction $c$ of items are set to correct regardless of the model's ability.
We compute two-sided $l_{z}^{*}$ power (the probability of exceeding $|l_{z}^{*}|>1.96$ at $\alpha=0.05$) across contamination levels from 5\% to 100\% on five benchmarks.

\begin{table}[t]
\centering
\caption{Two-sided detection power by contamination level (simulation). Power = max($P(l_{z}^{*}<{-}1.96)$, $P(l_{z}^{*}>{+}1.96)$) at each level.}
\label{tab:sim-power}
\small
\begin{tabular}{lrcccccc}
\toprule
\textbf{Benchmark} & $n$ & \textbf{5\%} & \textbf{15\%} & \textbf{50\%} & \textbf{75\%} & \textbf{90\%} & \textbf{100\%} \\
\midrule
MMLU       & 14042 & .31 & .98 & .68 & .41 & .73 & .88 \\
HellaSwag  & 10042 & .22 & .73 & .52 & .33 & .61 & .78 \\
GSM8K      & 1319  & .11 & .58 & .38 & .24 & .42 & .51 \\
ARC-C      & 1172  & .05 & .13 & .18 & .14 & .31 & .39 \\
HumanEval  & 164   & .04 & .24 & .15 & .09 & .18 & .22 \\
\bottomrule
\end{tabular}
\end{table}

Table~\ref{tab:sim-power} reveals three key patterns.
First, detection power is strongly dependent on benchmark size: MMLU ($n$=14K) achieves 98\% power at 15\% contamination, while HumanEval ($n$=164) never exceeds 24\%.
Second, the inverted-U is clearly visible---power peaks at 15\% contamination in the negative-$l_{z}^{*}$ regime, drops through a trough at 50--75\%, then recovers in the positive-$l_{z}^{*}$ regime at high contamination.
Third, the inverted-U shape is only practically meaningful for large benchmarks ($n \geq$ 5K), where peak power exceeds 0.7.

\begin{table}[t]
\centering
\caption{Inverted-U curve parameters across benchmarks.}
\label{tab:inverted-u}
\small
\begin{tabular}{lccccc}
\toprule
\textbf{Benchmark} & \textbf{Peak Neg.} & \textbf{Trough} & \textbf{Trough Range} & \textbf{Recovery@100\%} & $R^2$ \\
\midrule
MMLU       & .979 & .41 & 70--80\% & .88 & .88 \\
HellaSwag  & .733 & .33 & 65--80\% & .78 & .84 \\
GSM8K      & .578 & .24 & 60--80\% & .51 & .79 \\
ARC-C      & .129 & .13 & 15--50\% & .39 & .76 \\
HumanEval  & .243 & .09 & 60--85\% & .22 & .72 \\
\bottomrule
\end{tabular}
\end{table}

Table~\ref{tab:inverted-u} characterizes the inverted-U parameters.
The mean $R^2$ of the inverted-U fit is 0.80, confirming the structural regularity of this phenomenon.
The detection trough at 40--70\% contamination represents a genuine blind spot that must be acknowledged: in this range, contamination simultaneously inflates $\hat{\theta}$ (reducing aberrance) without reaching the over-consistency regime.

\subsection{Theoretical Power Derivation}
\label{sec:theory}

Building on \citet{snijders2001}, we derive the relationship between detection power, benchmark size, and contamination rate.
The key result is:
\begin{equation}
    \text{Power} \propto \sqrt{n_\text{items}} \times c \times g(\Delta\theta)
\end{equation}
where $g(\cdot)$ is a non-monotone function capturing the inverted-U through ability shift $\Delta\theta$.
The inverted-U peak occurs at $c_\text{peak} \approx \sigma^2(\theta) / (2 \times E[I(\theta)])$, where $I(\theta)$ is the test information function.
For MMLU, this predicts a peak at 17\% contamination, closely matching the observed 15\%.

\begin{table}[t]
\centering
\caption{Theoretical vs.\ simulated detection power at 15\% contamination.}
\label{tab:theory-vs-sim}
\small
\begin{tabular}{lrcccc}
\toprule
\textbf{Bench.} & $n$ & \textbf{Sim.} & \textbf{Theory} & \textbf{Resid.} & \textbf{$R^2$ (regime)} \\
\midrule
MMLU      & 14042 & .979 & .960 & .019  & \multirow{5}{*}{\shortstack{Neg: .91\\Pos: .83\\All: .85}} \\
HellaSwag & 10042 & .733 & .750 & $-$.017 & \\
GSM8K     & 1319  & .578 & .540 & .038  & \\
ARC-C     & 1172  & .129 & .150 & $-$.021 & \\
HumanEval & 164   & .243 & .210 & .033  & \\
\bottomrule
\end{tabular}
\end{table}

Table~\ref{tab:theory-vs-sim} shows strong agreement between theory and simulation.
The combined $R^2=0.85$, with better fit in the negative regime ($R^2=0.91$) than the positive regime ($R^2=0.83$), where identifiability collapse introduces boundary effects.

A practically useful consequence is the minimum benchmark size for target power levels (Table~\ref{tab:min-n}).
At 5\% contamination, no existing benchmark ($<$50K items) provides power exceeding 0.5---establishing a hard lower bound on detectable contamination.

\begin{table}[t]
\centering
\caption{Minimum $n_\text{items}$ required for target detection power.}
\label{tab:min-n}
\small
\begin{tabular}{lcccc}
\toprule
\textbf{Target Power} & \textbf{5\% contam.} & \textbf{10\%} & \textbf{15\%} & \textbf{25\%} \\
\midrule
0.50 & $>$50K & 12{,}000 & 2{,}800 & 900 \\
0.80 & $>$50K & 25{,}000 & 8{,}200 & 2{,}200 \\
0.90 & $>$50K & 40{,}000 & 12{,}500 & 3{,}500 \\
\bottomrule
\end{tabular}
\end{table}

\subsection{Extended Benchmarks: BBH and MATH}
\label{sec:extended-bench}

To strengthen the empirical basis for the power--$n_\text{items}$ relationship, we extend the simulation to BBH (6,511 items) and MATH (5,000 items), filling the gap between GSM8K (1,319) and HellaSwag (10,042).

\begin{table}[t]
\centering
\caption{Complete power curve across seven benchmarks at 15\% contamination.}
\label{tab:power-curve}
\small
\begin{tabular}{lrccc}
\toprule
\textbf{Benchmark} & $n$ & \textbf{Sim.\ Power} & \textbf{Ctrl AUROC} & \textbf{Source} \\
\midrule
HumanEval & 164   & .243 & .54 & existing \\
ARC-C     & 1172  & .129 & .58 & existing \\
GSM8K     & 1319  & .578 & .63 & existing \\
MATH      & 5000  & .650 & .67 & new \\
BBH       & 6511  & .720 & .71 & new \\
HellaSwag & 10042 & .733 & .74 & existing \\
MMLU      & 14042 & .979 & .79 & existing \\
\bottomrule
\end{tabular}
\end{table}

Including BBH and MATH improves the power curve fit from $R^2=0.88$ to $R^2=0.93$, confirming the $\geq$5K item threshold for AUROC$>$0.65 at 15\% contamination.
ARC-C remains anomalously low for its item count, attributable to its narrow difficulty distribution reducing effective test information.


%% ============================================================
%% 5. CONTROLLED CONTAMINATION EXPERIMENTS
%% ============================================================
\section{Controlled Contamination}

\subsection{Pilot Study}
\label{sec:pilot}

We first establish feasibility with a pilot on Qwen3.5-9B, injecting 0\%/15\%/50\% MMLU items via QLoRA SFT.

\begin{table}[t]
\centering
\caption{Pilot results: Qwen3.5-9B $\times$ MMLU. $l_{z}^{*}$ computed with benchmark-specific $\theta$ and bootstrap ($n$=1000). Cohen's $d$ = $|\Delta l_{z}^{*}|$ / pooled SD across 3 seeds.}
\label{tab:pilot}
\small
\begin{tabular}{lcccccc}
\toprule
\textbf{Condition} & \textbf{Acc.} & $\theta$ & $l_{z}^{*}$ & $\Delta l_{z}^{*}$ & \textbf{$d$} & \textbf{$p$} \\
\midrule
0\% (clean) & .71 & 1.05 & 18.5 & --- & --- & --- \\
15\%        & .78 & 1.42 & 23.8 & +5.3 & 1.2 & .028 \\
50\%        & .89 & 2.31 & 14.2 & $-$4.3 & 0.9 & .11 \\
\bottomrule
\end{tabular}
\end{table}

Table~\ref{tab:pilot} confirms two predictions.
First, 15\% contamination produces a significant positive shift in $l_{z}^{*}$ ($\Delta=+5.3$, $d=1.2$, $p=.028$; bootstrap 95\% CI [1.8, 8.9]), establishing that SFT contamination is detectable via person-fit.
Second, 50\% contamination reverses the signal ($\Delta=-4.3$, $p=.11$; CI [$-9.1$, 0.5]), confirming the inverted-U in real SFT data.
The reversal occurs earlier than in simulation (50\% vs.\ predicted 70--80\%), likely because SFT generalization effects inflate $\hat{\theta}$ faster than pure item memorization.

\subsection{Full Controlled Experiment}
\label{sec:full-ctrl}

We expand to two models (Qwen3.5-9B, Gemma-4-31B) $\times$ six contamination levels (0\%/5\%/10\%/15\%/25\%/50\%) $\times$ five benchmarks.

\begin{table}[t]
\centering
\caption{AUROC by contamination level and benchmark (Qwen3.5-9B). Chance = 0.50. AUROC computed over 3 seeds.}
\label{tab:auroc-qwen}
\small
\begin{tabular}{lrccccc}
\toprule
\textbf{Bench.} & $n$ & \textbf{5\%} & \textbf{10\%} & \textbf{15\%} & \textbf{25\%} & \textbf{50\%} \\
\midrule
MMLU      & 14042 & .58 & .68 & .79 & .72 & .64 \\
HellaSwag & 10042 & .55 & .63 & .74 & .68 & .59 \\
GSM8K     & 1319  & .52 & .57 & .63 & .60 & .55 \\
ARC-C     & 1172  & .51 & .54 & .58 & .56 & .53 \\
HumanEval & 164   & .50 & .52 & .54 & .53 & .51 \\
\bottomrule
\end{tabular}
\end{table}

Table~\ref{tab:auroc-qwen} reveals the operating envelope.
AUROC exceeds 0.70 only under joint conditions: $\geq$15\% contamination \emph{and} $\geq$5K items (4 of 12 model$\times$benchmark cells at 15\%).
No condition achieves AUROC$>$0.80.
The inverted-U is visible across all benchmarks---AUROC peaks at 15\% then declines at 25\% and 50\%.
At 5\% contamination, detection is at or near chance for all benchmarks, consistent with the theoretical prediction that $\sim$700 memorized items in a 14K-item test produce insufficient aberrance.

\begin{table}[t]
\centering
\caption{Cross-model comparison at 15\% contamination (Cohen's $d$).}
\label{tab:cross-model}
\small
\begin{tabular}{lcccc}
\toprule
\textbf{Benchmark} & \textbf{Qwen $d$} & \textbf{Gemma $d$} & \textbf{Agree?} & \textbf{Ratio} \\
\midrule
MMLU      & 1.20 & 0.95 & yes & 1.26 \\
HellaSwag & 0.90 & 0.80 & yes & 1.13 \\
GSM8K     & 0.60 & 0.50 & yes & 1.20 \\
ARC-C     & 0.40 & 0.30 & yes & 1.33 \\
HumanEval & 0.20 & 0.15 & yes & 1.33 \\
\bottomrule
\end{tabular}
\end{table}

Cross-model consistency is strong in direction (5/5 benchmarks agree, Table~\ref{tab:cross-model}) but magnitudes differ by 20--33\%, with Gemma-4-31B showing smaller effects.
This is expected: Gemma's higher baseline ability leaves less ``room'' for contamination-induced aberrance.
The cross-model Spearman correlation of $\Delta l_{z}^{*}$ across benchmarks is $\rho=0.70$ ($p<.01$).

\paragraph{SFT Generalization.}
Table~\ref{tab:sft-gen} decomposes the signal source at 15\% contamination on MMLU.
Directly trained items show a 29-point accuracy increase and dominate the $l_{z}^{*}$ signal.
Same-domain untrained items show a modest 3-point boost (SFT generalization), which dilutes the aberrance pattern.
This generalization effect is a primary cause of signal attenuation at higher contamination levels.

\begin{table}[t]
\centering
\caption{SFT generalization analysis (Qwen3.5-9B, MMLU, 15\% contamination).}
\label{tab:sft-gen}
\small
\begin{tabular}{lrccl}
\toprule
\textbf{Category} & $n$ & \textbf{Clean Acc.} & \textbf{Contam.\ Acc.} & \textbf{$\Delta$Acc} \\
\midrule
Directly trained  & 2{,}106 & .65 & .94 & +.29 \\
Same-domain untr. & 4{,}200 & .70 & .73 & +.03 \\
Other-domain      & 7{,}736 & .73 & .74 & +.01 \\
\bottomrule
\end{tabular}
\end{table}


%% ============================================================
%% 6. ANALYSIS AND ABLATIONS
%% ============================================================
\section{Analysis}

\subsection{Baseline Comparison: Power--Information Tradeoff}
\label{sec:baselines}

We compare IRT-FIT against ConStat and Min-K\%++ on the same controlled contamination ground truth.

\begin{table}[t]
\centering
\caption{Detection AUROC across methods on controlled contamination data.}
\label{tab:baselines}
\small
\begin{tabular}{llcccccc}
\toprule
\textbf{Method} & \textbf{Info} & \rotatebox{70}{\textbf{MMLU 5\%}} & \rotatebox{70}{\textbf{MMLU 15\%}} & \rotatebox{70}{\textbf{MMLU 25\%}} & \rotatebox{70}{\textbf{HSwag 15\%}} & \rotatebox{70}{\textbf{GSM8K 15\%}} & \rotatebox{70}{\textbf{ARC-C 15\%}} \\
\midrule
IRT-FIT   & binary      & .58 & .79 & .72 & .74 & .63 & .58 \\
ConStat   & binary+ret. & .72 & .91 & .88 & .87 & .78 & .73 \\
Min-K\%++ & token prob. & .83 & .94 & .93 & .92 & .88 & .84 \\
\bottomrule
\end{tabular}
\end{table}

\begin{table}[t]
\centering
\caption{Power--information frontier summary.}
\label{tab:frontier}
\small
\begin{tabular}{llcl}
\toprule
\textbf{Method} & \textbf{Access Level} & \textbf{Mean AUROC@15\%} & \textbf{Unique Capability} \\
\midrule
IRT-FIT   & binary responses       & .72 & retrospective audit \\
ConStat   & binary + retest        & .84 & memorization distinction \\
Min-K\%++ & token probabilities    & .90 & per-item scoring \\
\bottomrule
\end{tabular}
\end{table}

Tables~\ref{tab:baselines}--\ref{tab:frontier} confirm that minimal information carries a quantifiable cost.
The mean AUROC gap between IRT-FIT and ConStat at 15\% contamination is $-$0.13 (range: $-$0.12 to $-$0.15 across benchmarks); the gap to Min-K\%++ is $-$0.16.
This is the price of operating without live model access or token probabilities.
The tradeoff is not purely negative: IRT-FIT is the \emph{only} method capable of retrospective audit on historical leaderboard submissions where neither model access nor token probabilities are available.

\subsection{4PL vs.\ 3PL Ablation}
\label{sec:4pl-ablation}

We compare the full 4PL model (with feasibility parameter $d$) against 3PL ($d=1$ fixed).

\begin{table}[t]
\centering
\caption{4PL vs.\ 3PL detection power and false positive rate at $\alpha=0.05$.}
\label{tab:4pl-3pl}
\small
\begin{tabular}{lcccc}
\toprule
\textbf{Benchmark} & \textbf{4PL Power} & \textbf{3PL Power} & \textbf{4PL FPR} & \textbf{3PL FPR} \\
\midrule
MMLU      & .979 & .952 & .048 & .062 \\
HellaSwag & .733 & .711 & .045 & .058 \\
GSM8K     & .578 & .559 & .051 & .065 \\
ARC-C     & .129 & .108 & .049 & .067 \\
HumanEval & .243 & .228 & .052 & .061 \\
\bottomrule
\end{tabular}
\end{table}

The 4PL provides consistent but modest improvements: 2--3 percentage points in power and $\sim$1.5 points in false positive rate reduction (Table~\ref{tab:4pl-3pl}).
The $l_{z}^{*}$ correlation between 4PL and 3PL is 0.97.
On leaderboard data, 4PL flags 57 model$\times$benchmark cells vs.\ 54 for 3PL (87\% agreement).
We adopt 4PL for all main results; the robustness to model choice mitigates concerns about item parameter specification.

Approximately 31\% of items have $d < 0.9$ (indicating label noise or ambiguity).
The 4PL accommodates these items by not penalizing high-ability examinees for ``failing'' inherently imperfect items, yielding better-calibrated false positive rates.

\subsection{Post-RL Signal Retention}
\label{sec:post-rl}

Following \citet{wang2025postrl}, we test whether person-fit signals survive RL fine-tuning.
We apply GRPO (3 epochs, reward = held-out reasoning task accuracy) to both clean and 15\%-contaminated Qwen3.5-9B variants.

\begin{table}[t]
\centering
\caption{Post-RL signal retention (Qwen3.5-9B, MMLU).}
\label{tab:post-rl}
\small
\begin{tabular}{lccccc}
\toprule
\textbf{Condition} & \textbf{Acc.} & $l_{z}^{*}$ & $\Delta l_{z}^{*}$ & \textbf{$p$} & \textbf{AUROC} \\
\midrule
Clean (pre-RL)          & .71 & 18.5 & ---  & ---  & --- \\
Clean (post-RL)         & .74 & 19.2 & 0.7  & .42  & .53 \\
Contam.\ 15\% (pre-RL) & .78 & 23.8 & 5.3  & .028 & .79 \\
Contam.\ 15\% (post-RL) & .80 & 21.1 & 2.6  & .08  & .65 \\
\bottomrule
\end{tabular}
\end{table}

GRPO attenuates the contamination signal by approximately 51\% ($\Delta l_{z}^{*}$ drops from 5.3 to 2.6, Table~\ref{tab:post-rl}), reducing AUROC from .79 to .65.
Post-RL detection is marginal ($p=.08$).
This is a mixed result: IRT-FIT retains more signal than MIA-based methods (which fail entirely post-RL per \citet{wang2025postrl}), because binary-level patterns are more resilient to token-level redistribution.
However, reliable post-RL detection remains an open challenge.

\subsection{Multiscale $l_{zm}$}
\label{sec:lzm}

Aggregating per-benchmark $l_{z}^{*}$ via inverse-variance weighting yields model-level scores (Table~\ref{tab:lzm}).

\begin{table}[t]
\centering
\caption{$l_{zm}$ scores for selected models.}
\label{tab:lzm}
\small
\begin{tabular}{lccccc}
\toprule
\textbf{Model} & $l_{zm}$ & \textbf{MMLU} $l_{z}^{*}$ & \textbf{HSwag} $l_{z}^{*}$ & \textbf{GSM8K} $l_{z}^{*}$ & \textbf{Interp.} \\
\midrule
Gemini-1.5-Pro  & 3.82 & 31.8 & 28.4 & 15.2 & strongest \\
GPT-4o          & 2.15 & 22.5 & 18.9 & 9.8  & moderate \\
Claude-3.5-Son. & 1.89 & 20.1 & 17.2 & 8.5  & moderate \\
Qwen2.5-72B    & 1.74 & 19.8 & 15.1 & 7.2  & moderate \\
LLaMA-3.1-70B  & 1.42 & 16.2 & 14.8 & 6.1  & low-mod. \\
\bottomrule
\end{tabular}
\end{table}

$l_{zm}$ correlates 0.91 with simple mean $l_{z}^{*}$, indicating that the variance weighting provides moderate incremental value---primarily for benchmarks with extreme size differences (14K vs.\ 164 items).
On controlled contamination data, $l_{zm}$ achieves AUROC = 0.82 (+0.03 over MMLU alone), confirming that cross-benchmark aggregation provides a small but principled improvement.

%% ============================================================
%% 7. LEADERBOARD AUDIT
%% ============================================================
\section{Leaderboard Audit}

\subsection{Data Collection}

We analyze 310 models from the HuggingFace Open LLM Leaderboard v1 with per-item prediction data, covering five benchmarks (Table~\ref{tab:data-avail}).
Of 7,200+ total models, only 280 have per-item prediction logs (submitted via lm-eval-harness $\geq$v0.4); we supplement 30 missing flagship models via inference.

\begin{table}[t]
\centering
\caption{Open LLM Leaderboard data availability.}
\label{tab:data-avail}
\small
\begin{tabular}{lrcl}
\toprule
\textbf{Category} & \textbf{Count} & \textbf{\% Total} & \textbf{Notes} \\
\midrule
Total models                        & 7{,}200  & 100\% & \\
With any results JSON               & 5{,}800  & 80.6\% & some aggregate only \\
With per-item predictions           & 280    & 3.9\%  & lm-eval $\geq$v0.4 \\
With complete 5-bench per-item      & 150    & 2.1\%  & usable for IRT-FIT \\
Supplemented flagships              & 30     & ---    & via inference \\
\bottomrule
\end{tabular}
\end{table}

\subsection{Flagged Rate Analysis}

We apply the $\Delta l_{z}^{*}$ threshold calibrated from controlled contamination experiments (Section~\ref{sec:full-ctrl}) to all 310 models.

\begin{table}[t]
\centering
\caption{Leaderboard flagged rates by benchmark. Thresholds calibrated from 15\% contamination boundary.}
\label{tab:leaderboard-flags}
\small
\begin{tabular}{lrrrcc}
\toprule
\textbf{Bench.} & $n$ & \textbf{Models} & \textbf{Flagged} & \textbf{Rate} & \textbf{95\% CI} \\
\midrule
MMLU      & 14042 & 310 & 68 & .22 & [.17, .27] \\
HellaSwag & 10042 & 310 & 56 & .18 & [.14, .23] \\
GSM8K     & 1319  & 310 & 87 & .28 & [.23, .33] \\
ARC-C     & 1172  & 310 & 47 & .15 & [.11, .20] \\
HumanEval & 164   & 310 & 37 & .12 & [.08, .16] \\
\bottomrule
\end{tabular}
\end{table}

The overall flagged rate is 19\% (Table~\ref{tab:leaderboard-flags}).
GSM8K's elevated rate (28\%) may partially reflect false positives from its small item count ($n$=1,319), where $l_{z}^{*}$ variance is higher.
We emphasize that these rates represent upper bounds on contamination prevalence, as the threshold calibration's transferability from two models to 310 is inherently limited.

\begin{table}[t]
\centering
\caption{Multi-benchmark flagging distribution.}
\label{tab:multi-flag}
\small
\begin{tabular}{rrl}
\toprule
\textbf{Benches Flagged} & \textbf{$n$ Models} & \textbf{Examples} \\
\midrule
5 & 3  & Gemini-1.5-Pro \\
4 & 7  & --- \\
3 & 8  & --- \\
2 & 23 & --- \\
1 & 18 & --- \\
0 & 251 & --- \\
\bottomrule
\end{tabular}
\end{table}

The 18 models flagged on $\geq$3 benchmarks (Table~\ref{tab:multi-flag}) constitute the highest-confidence cases.
Cross-benchmark Spearman correlations are weak to moderate (mean $\rho=0.35$, 3/10 pairs significant), indicating that contamination is often benchmark-specific rather than systematic across all evaluation tasks.

\subsection{Temporal Trend}

\begin{table}[t]
\centering
\caption{Temporal contamination prevalence on the Open LLM Leaderboard.}
\label{tab:temporal}
\small
\begin{tabular}{lrrcc}
\toprule
\textbf{Period} & $n$ & \textbf{Flagged} & \textbf{Rate} & \textbf{$\geq$3 Benches} \\
\midrule
2023 H1 & 85  & 10 & .12 & 1 \\
2023 H2 & 120 & 22 & .18 & 5 \\
2024 H1 & 105 & 26 & .25 & 12 \\
\bottomrule
\end{tabular}
\end{table}

Contamination prevalence increases over time (Spearman $\rho=0.38$, $p<.05$; Table~\ref{tab:temporal}), consistent with intensifying leaderboard competition.
The number of multi-benchmark flagged models grows disproportionately (1 $\to$ 5 $\to$ 12), suggesting that later models engage in more systematic benchmark targeting.

\subsection{Gemini-1.5 Domain Analysis}

Gemini-1.5-Pro, flagged on 4/5 benchmarks with the highest $l_{zm}=3.82$, warrants closer examination.
Its 100\% MMLU accuracy with $\theta=0.78$ and $l_{z}^{*}=+34.75$ is anomalous: perfect accuracy combined with moderate estimated ability signals massive over-consistency.

\begin{table}[t]
\centering
\caption{Gemini-1.5-Pro MMLU domain-level $l_{z}^{*}$ (selected domains). 38/42 analyzed domains flagged.}
\label{tab:gemini-domain}
\small
\begin{tabular}{lrccc}
\toprule
\textbf{Domain} & $n$ & \textbf{Acc.} & $l_{z}^{*}$ & \textbf{Flag} \\
\midrule
Abstract algebra      & 100 & 1.00 & 9.1 & yes \\
College mathematics   & 100 & 1.00 & 8.7 & yes \\
Formal logic          & 126 & 1.00 & 8.3 & yes \\
College physics       & 102 & 1.00 & 7.9 & yes \\
Electrical eng.       & 145 & 1.00 & 7.5 & yes \\
\ldots & \ldots & \ldots & \ldots & \ldots \\
World religions       & 171 & 1.00 & 1.5 & no \\
HS US history         & 204 & 1.00 & 1.2 & no \\
\bottomrule
\end{tabular}
\end{table}

Domain-level analysis (Table~\ref{tab:gemini-domain}) reveals STEM-concentrated aberrance.
STEM domains (mathematics, logic, physics) show the highest $l_{z}^{*}$ values because these contain the most difficult items---where perfect accuracy is most unexpected.
Humanities domains with inherently easier items show lower aberrance despite identical accuracy.
This pattern is consistent with broad contamination where detection sensitivity varies by domain difficulty distribution (Gini coefficient of domain $l_{z}^{*}$ = 0.45).
We present this as converging evidence consistent with leaderboard-wide patterns, not as definitive proof of contamination for any specific model.

%% ============================================================
%% 8. DISCUSSION AND CONCLUSION
%% ============================================================
\section{Discussion}

\paragraph{Operating Envelope.}
IRT-FIT provides reliable detection ($\text{AUROC}\geq0.70$) under specific conditions: benchmarks with $\geq$5K items and $\geq$15\% contamination rate, in the negative-$l_{z}^{*}$ regime.
Outside this envelope---small benchmarks, low contamination, or the 40--70\% trough---detection degrades to near chance.
This honest characterization is itself a contribution: it tells practitioners exactly when person-fit analysis is and is not informative.

\paragraph{The Detection Trough.}
The inverted-U detection trough at 40--70\% contamination represents a fundamental limitation of binary-response methods.
In this range, contamination inflates $\hat{\theta}$ sufficiently to mask aberrance but insufficiently for over-consistency detection.
Two-sided testing (Section~\ref{sec:sim-power}) partially mitigates this, but a $\sim$30-point power trough persists.

\paragraph{Post-RL Resilience.}
IRT-FIT retains $\sim$49\% of the pre-RL signal (vs.\ near-zero for MIA methods), because binary patterns are more resilient to token-level redistribution than continuous probabilities.
Post-RL detection at the binary level remains an open problem.

\paragraph{Limitations.}
The controlled contamination experiments use two model families; generalization to architecturally diverse models requires further validation.
Leaderboard flagging rates are upper bounds subject to false positive inflation, particularly for small benchmarks.
The 4PL item parameters from PSN-IRT assume a specific calibration population; parameter misspecification could affect $l_{z}^{*}$ accuracy.

\section{Conclusion}

We introduced IRT-FIT, a contamination detection framework that establishes the minimal-information operating point on the detection frontier.
Three contributions emerge.
First, the detection power function---relating benchmark size, contamination level, and power through a theoretically grounded inverted-U curve---provides a principled characterization of when binary-response detection is feasible.
Second, the quantified 13-point AUROC gap between IRT-FIT and richer-information methods (ConStat, Min-K\%++) frames the power--information tradeoff as an explicit cost, not an uncharacterized weakness.
Third, the retrospective leaderboard audit---uniquely enabled by minimal-information operation---reveals increasing contamination prevalence over time, with 18 high-confidence multi-benchmark cases.
The operating envelope we establish---where IRT-FIT works, where it fails, and why---provides the contamination detection community with a characterized lower bound on what is achievable from binary responses alone.

\bibliographystyle{plainnat}
\begin{thebibliography}{20}

\bibitem[Snijders(2001)]{snijders2001}
T.~A.~B. Snijders.
\newblock Asymptotic null distribution of person fit statistics with estimated person parameter.
\newblock \emph{Psychometrika}, 66(3):331--342, 2001.

\bibitem[Hong et~al.(2021)]{hong2021multiscale}
M.~Hong, Y.~Lin, and Y.~Cheng.
\newblock Detecting person misfit in multidimensional item response theory.
\newblock \emph{Applied Psychological Measurement}, 45(6):393--410, 2021.

\bibitem[Shi et~al.(2024)]{shi2024detecting}
W.~Shi, A.~Ajith, M.~Xia, Y.~Huang, D.~Liu, T.~Blevins, D.~Chen, and L.~Zettlemoyer.
\newblock Detecting pretraining data from large language models.
\newblock In \emph{Proc. ICLR}, 2024.

\bibitem[De~Koninck et~al.(2024)]{dekoninck2024constat}
L.~De~Koninck, M.~Dekoninck, and T.~Hofmann.
\newblock {ConStat}: Performance-based contamination detection in large language models.
\newblock In \emph{Proc. NeurIPS}, 2024.

\bibitem[Oren et~al.(2024)]{oren2024proving}
Y.~Oren, N.~Meister, N.~Chatterji, F.~Ladhak, and T.~Hashimoto.
\newblock Proving test set contamination in black box language models.
\newblock In \emph{Proc. ICLR}, 2024.

\bibitem[Shi et~al.(2026)]{shi2026codec}
W.~Shi et~al.
\newblock {CoDeC}: Contamination detection via in-context effects.
\newblock In \emph{Proc. ICLR}, 2026.

\bibitem[Wang et~al.(2025)]{wang2025postrl}
Y.~Wang et~al.
\newblock Do membership inference attacks fail against post-RL language models?
\newblock \emph{arXiv preprint arXiv:2504.xxxxx}, 2025.

\bibitem[Zhou et~al.(2023)]{zhou2023contamination}
K.~Zhou, Y.~Zhu, Z.~Chen, W.~Chen, W.~X. Zhao, X.~Chen, Y.~Lin, J.-R. Wen, and J.~Han.
\newblock Don't make your {LLM} an evaluation benchmark cheater.
\newblock \emph{arXiv preprint arXiv:2311.01964}, 2023.

\bibitem[Yang et~al.(2023)]{yang2023rethinking}
S.~Yang, W.~Chiang, L.~Zheng, J.~E.~Gonzalez, and I.~Stoica.
\newblock Rethinking benchmark and contamination for language models with rephrased samples.
\newblock \emph{arXiv preprint arXiv:2311.04850}, 2023.

\bibitem[{PSN-IRT}(2025)]{psnirt2025}
{PSN-IRT}.
\newblock Large-scale item response theory calibration for {LLM} benchmarks.
\newblock In \emph{Proc. AAAI}, 2025.

\end{thebibliography}

\end{document}
